% !PRINTED 330
\weberpage{330}

\begin{center}
\gesperrt{Tenth Chapter.}\\[4pt]
\textbf{Algebraic Numbers and Forms.}
\end{center}

\begin{center}
\rule{0.25\textwidth}{0.4pt}
\end{center}

\begin{center}
\S\,87. \gesperrt{Ideals and Forms in Algebraic Fields.}
\end{center}

The interest which the theory of elliptic functions has for the algebraist springs from its application to the theory of quadratic irrational numbers, to which these functions stand in a relation analogous to that of the roots of unity to the rational numbers. They offer us the first and hitherto only example, more exactly known in its laws, of a region of algebraic numbers going beyond the cyclotomic numbers. In order to be able to present the theory of these numbers more fully, we must first premise some theorems from the theory of forms and of algebraic fields, and especially of quadratic forms.

In Bd. II, \S\,163, 169 a connection was found between the ideals of an algebraic field $\Omega$ of the $n$th degree and certain homogeneous functions of the $n$th degree of $n$ variables. Before we apply this to quadratic fields, let the general theorems be briefly recalled.

Let $\mathfrak{a}$ be an ideal of a field $\Omega$ and let $\alpha_1, \alpha_2, \ldots, \alpha_n$ be a basis of this ideal. Further, let $\omega_1, \omega_2, \ldots, \omega_n$ be a minimal basis of $\Omega$. Then
\weq{1}{(\alpha_1, \alpha_2, \ldots, \alpha_n) = (A)(\omega_1, \omega_2, \ldots, \omega_n),}
if $(A)$ denotes a \gesperrt{linear substitution with integral coefficients}. In order to pass over to another basis $\alpha_1', \alpha_2', \ldots, \alpha_n'$ of $\mathfrak{a}$, one makes a linear substitution $L$ with the determinant $\pm 1$,
\weq{2}{(\alpha_1', \alpha_2', \ldots, \alpha_n') = (L)(\alpha_1, \alpha_2, \ldots, \alpha_n),}
from which there results
\weq{3}{(\alpha_1', \alpha_2', \ldots, \alpha_n') = (L)(A)(\omega_1, \omega, \ldots, \omega_n)^{\text{1)}}.} % [sic] printed \textquotedbl{}$\omega$\textquotedbl{} without subscript 2 in the second entry

\medskip
\noindent\rule{0.3\textwidth}{0.4pt}\\
$^{1}$) On linear substitutions and their composition cf. the sixth chapter of the second volume.
